\subsection*{36. A Beachhead Must Fit the Founders' Delivery Competence}

A small company can be an expert in the problem and inexperienced in delivering the promised solution at scale. An engineer may know model evaluation but not educational assessment. A teacher may be excellent live and unable to coordinate 30 parallel learners' projects. A researcher may understand the science but be unready to manage enterprise information security.

Evaluate the full service system: acquisition, onboarding, delivery, support, error recovery, payment, and measurable outcomes. The beachhead should let the company demonstrate competence while developing adjacent abilities gradually. A market requiring every missing capability at once is a poor first choice even if it is attractive later.

\subsection*{37. The First Proof Should Survive Without the Founder's Charisma}

Founder-driven trust can produce the first transactions. But ask what the evidence would look like to a buyer who has never met the founder. Do you have an inspectable work sample, a rigorous method, a paid pilot result, documented feedback, a reference permitted to speak, or an independently assessable artifact?

Choosing a beachhead with \textbf{portable proof} matters. If customers' results are always confidential or deeply subjective, the business may need a sophisticated reference process. If a technical institute can produce consented, independently reviewable learner artifacts, it may have a clearer way to demonstrate progress.

Proof must correspond to the claim. A founder's impressive demo is evidence of founder competence, not an institutional record. An alumni artifact is evidence of that learner's work, not automatically causal proof of employment impact.

\subsection*{38. Use a Hard-Gate Method Before Any Ranking}

A weighted score cannot repair a fatal feasibility problem. Before comparing opportunities, ask whether each candidate satisfies essential conditions:

\begin{table}[htbp]
\centering
\small
\begin{tabularx}{\columnwidth}{l X}
\toprule
\textbf{Hard Gate} & \textbf{Failure Means} \\
\midrule
Legitimate to offer & Violates a legal, ethical, employment, or mission boundary \\
Plausible paid demand & No identifiable buyer or feasible purchasing mechanism exists \\
Reachable buyer & Founder has no realistic, affordable route to even a few prospects \\
Deliverable promise & Cannot honestly perform the job with current or obtainable resources \\
Survivable downside & Pilot creates liabilities or costs the company cannot absorb \\
Testable in runway & Evidence will arrive only after resources are exhausted \\
\bottomrule
\end{tabularx}
\caption{The 6 non-negotiable feasibility gates.}
\label{tab:ch7-hard-gates}
\end{table}

A failed gate need not permanently eliminate the opportunity. It can mean \emph{defer until a partner, license, instructor, channel, or reserve is acquired}. Do not bury the failure inside an average score.

\subsection*{39. Score Uncertainty, Not Just Enthusiasm}

For candidates that pass the gates, use a comparison table. Suggested dimensions include urgency; reachable economic buyers; prior purchasing behavior; contrast with existing alternatives; founder competence; minimum viable delivery cost; time to decision; customer concentration risk; reputation transfer; asset compounding; and adjacency potential.

Scores are useful to expose disagreements. They are not scientific forecasts. Assign \textbf{evidence confidence separately}: directly observed, strongly inferred, weakly inferred, or unknown. A score of 5 based on imagination should not outrank a score of 3 supported by completed purchases.

A practical scorecard might rate strength 0 to 3 and evidence 0 to 3, but the numbers are merely labels for discussion. If desired, include a conservative evidence-adjusted ranking for screening; then inspect the raw claims. Never use multiplication to launder guesses into objective probabilities.

\subsection*{40. The Most Important Question in the Comparison Is Often Missing}

Founders ask, ``How big is the opportunity?'' They should also ask: \textbf{What is the next claim we must establish for this opportunity to become investable?}

For a cohort, perhaps the missing claim is that capable professionals will pay for expert review. For a university partnership, perhaps it is that the dean can authorize a funded pilot within the semester. For enterprise training, perhaps it is that an engineering leader has a business performance problem and budget rather than an interest in an impressive workshop.

The best next beachhead experiment attacks the weakest decisive link. If the offer cannot be delivered, experiment with delivery before buying traffic. If no payer is identifiable, investigate purchasing before refining curriculum. If the product works but the segment is impossible to reach, test distribution, not more product features.

\subsection*{41. 3 Candidates and a Falsifiable Selection Thesis}

Imagine a hypothetical founder choosing among 3 segments:

\begin{table*}[t]
\centering
\small
\begin{tabularx}{\textwidth}{>{\raggedright\arraybackslash}p{3.2cm} >{\raggedright\arraybackslash}p{3.0cm} >{\raggedright\arraybackslash}p{2.8cm} >{\raggedright\arraybackslash}p{3.2cm} >{\raggedright\arraybackslash}X}
\toprule
\textbf{Candidate} & \textbf{Need} & \textbf{Access} & \textbf{Distinctive Possibility} & \textbf{Decisive Uncertainty} \\
\midrule
Early-career beginners seeking broad AI skills & Widespread curiosity and career anxiety & Large but competitive paid channels & Founder technical expertise & Can an unknown provider win economical students without placement claims? \\
Working developers responsible for live AI & Potentially urgent implementation problems & Practitioner communities, peer introductions & Deep technical review and adult teaching & Will they pay and sustain structured learning, or prefer ad hoc support? \\
Corporate teams beginning applied AI & Observable workflow and capability gaps & Targeted managers; corporate procurement & Training tied to system deliverables & Can a small new provider access budgets and prove business impact? \\
\bottomrule
\end{tabularx}
\caption{Screening 3 hypothetical entry segments on decisive uncertainties.}
\label{tab:ch7-three-candidates}
\end{table*}

The founder should not declare the second group victorious because its narrative is appealing. Instead, write a hypothesis: ``We favor this group because our expertise makes its problem legible, a complete proof experience is feasible, and we can reach at least a few eligible buyers. We will change direction if paid demand, learner time, or economic delivery fails predefined thresholds.''

This preserves decisiveness without pretending the future is known.

\subsection*{42. Selection Should Seek Asymmetric Information Gain}

An experiment is unusually attractive when it is cheap to run, its downside is bounded, and its result informs multiple important decisions. A short paid engineering review might test buyer urgency, ability to pay, credibility, teaching style, faculty workload, and prospective curriculum design. A 6-month program tests those too, but after making a far larger commitment.

This is an application of Chapter 4's value-of-information reasoning: \textbf{prefer experiments that eliminate consequential uncertainty before irreversible spending.} A low-cost test that reveals a fatal assumption is successful research even when it produces little revenue.

However, free curiosity is not paid demand. A workshop's usefulness does not automatically validate a months-long course. Specify which inference the test can and cannot justify.

\subsection*{43. Select 1 Operational Beachhead; Preserve a Watchlist}

After sufficient comparison, name a primary beachhead. Keep the other opportunities in an explicit parking lot with conditions for reconsideration. A watchlist is not a concealed second product pipeline.

For each parked candidate, record what evidence would make it outrank the current choice: a signed buyer request, a credible partner offering funded distribution, a new regulation that creates urgent demand, or a successful test showing drastically better economics.

The organization can maintain strategic optionality by collecting low-cost intelligence without investing equally in every possibility. This is different from indecision. \textbf{Options are valuable when they preserve choices; they become expensive when exercised prematurely.}
